% Book pp. 53-57 (PDF pp. 54-58)
\subsection{Sum of an Arithmetic Progression (AP)}

\begin{activity}{Activity 2.2}
\begin{enumerate}
  \item Working in pairs or small groups, work through the following steps to
  discover how to find the general sum of an arithmetic sequence, ie a
  sequence with a common difference.
  \item Let $Sn$ denote the sum of the first $n$ terms of a linear sequence, where
  ``$a$'' is the first term and ``$d$'' the common difference.
  \item $Sn = a + (a+d) + (a+2d) + \ldots + [a+(n-1)d] \ \ldots\ldots\ldots\ldots\ldots\ (1)$
  \item Writing the terms in the reverse order, we get :
  \item $Sn = [a+(n-1)d] + [a+2d(n-2)] + \ldots + (a+2d) + (a+d) + a \ \ldots\ldots\ldots\ (2)$

  Adding (1) and (2), we get
  \[
    2Sn = [2a+(n-1)d] + [2a+(n-1)d] + [2a+(n-1)] + \ldots = n[2a+(n-1)d]
  \]
  (Since we are adding $n$ terms we have $n$ lots of this expression.)
  \[
    \therefore Sn = \frac{n}{2}\big[2a + (n-1)d\big]
  \]
  \item Let $U_1 = a$,
  \item Then, $Sn = \frac{n}{2}\big[2U_1 + (n-1)d\big]$
  \item Similarly, $Sn = \frac{n}{2}\big[U_1 + \{U_1 + (n-1)d\}\big]$
  \item But $U_n = U_1 + (n-1)d$
  \item By substitution, we get
  \item $S_n = \frac{n}{2}\big[U_1 + U_n\big]$ (Remember that $U_1$ = the first term and $U_n$ = the last term)
\end{enumerate}
\booknote{the book prints the second term of line (2) as $[a+2d(n-2)]$ and the third bracket of
$2Sn$ as $[2a+(n-1)]$; the intended expressions are $[a+(n-2)d]$ and $[2a+(n-1)d]$.}
\end{activity}

\begin{example}{Example 2.2}
\begin{enumerate}
  \item[a.] Find the sum of the first 35 terms of the linear sequence 2, 5, 8\ldots
\end{enumerate}

\solution
\begin{align*}
  a &= 2,\ d = 3 \text{ and } n = 35 \\
  Sn &= \frac{n}{2}\big[2a + (n-1)d\big] \\
  S_{35} &= \frac{35}{2}\big[2(2) + (35-1)(3)\big] = 1855
\end{align*}

\begin{enumerate}
  \item[b.] A sequence is defined by $U_{n+1} = 3 + U_n$ and $U_1 = 1$
  \item[c.] Find the sum of the first ten term of the sequence.
\end{enumerate}

\solution
\begin{align*}
  U_{n+1} &= 3 + U_n \\
  U_1 &= 1\ (= a) \\
  U_2 &= 4
\end{align*}
Common difference $d = u_2 - u_1 = 4 - 1 = 3$
\[
  S_n = \frac{n}{2}[2a + (n-1)d],
\]
where $a$ = first term, $d$ the common difference and $n$ the number of terms
\begin{align*}
  S_{10} &= \frac{10}{2}[2(1) + (10-1)3] \\
  S_{10} &= 5[2 + 9(3)] \\
  S_{10} &= 5[29] \\
  S_{10} &= 145
\end{align*}

\begin{enumerate}
  \item[d.] Find the the sum of the first 10 positive even numbers:
\end{enumerate}

\solution
\begin{align*}
  S_{10} &= 2 + 4 + 6 + 8 + 10 \cdots + 18 + 20 \\
  S_n &= \frac{n}{2}[U_1 + U_n] = \frac{n}{2}[\textit{first term} + \textit{last term}] \\
  S_{10} &= \frac{10}{2}[2 + 20] = 5 \times 22 = 110
\end{align*}
\end{example}

\begin{example}{Example 2.3}
A conference hall has 189 seats, arranged in 9 rows. The organiser wants to
arrange the seats such that the difference in the number of seats between any two
consecutive rows is 5 seats. If there are 9 rows in total, how many seats should be
in the third row?

\solution
Number of rows $n = 9$

Capacity of the room $S_9 = 189$

Difference $d = 5$

Let the number of seats in the first row $= a$
\begin{align*}
  S_n &= \frac{n}{2}[2a + (n-1)d] \\
  S_9 &= \frac{9}{2}[2a + (9-1)5] \\
  189 &= \frac{9}{2}[2a + (8)5] \\
  378 &= 9[2a + (8)5] \\
  378 &= 18a + 360 \\
  18a &= 378 - 360 \\
  18a &= 18 \\
  a &= 1
\end{align*}
Number of seats in the third row $a + 2d \Rightarrow 1 + 2(5) = 11$
\end{example}

\subsection{Sum of a Geometric Progression (GP)}

Remember that a geometric progression has the first term, $a$ and a common ratio, $r$.

The position of terms represented by $n$ can be written as:
\begin{align*}
  S_n &= a + ar + ar^2 + ar^3 + \ldots + ar^{n-2} + ar^{n-1} \ \ldots\ldots\ (1) \\
  rS_n &= ar + ar^2 + ar^3 + \ldots\ldots + ar^{n-1} + ar^{n} \ \ldots\ldots\ (2)
\end{align*}
$\textit{eqn } (2) - \textit{eqn } (1)$
\begin{align*}
  rSn - Sn &= -a + ar^n \\
  S_n(r-1) &= ar^n - a \\
  S_n &= \frac{a(r^n - 1)}{r - 1} = \frac{-a(1 - r^n)}{-(1 - r)} \\
  S_n &= \frac{a(r^n - 1)}{r - 1},\ r \ne 1 \\
  S_n &= \frac{a(1 - r^n)}{(1 - r)},\ r \ne 1
\end{align*}

\begin{example}{Example 2.4}
Find the sum of the first seven terms of the sequence 2, 8, 32\ldots

\solution
\begin{align*}
  S_n &= \frac{a(r^n - 1)}{r - 1},\ r \ne 1 \\
  a &= 2,\ r = 4 \text{ and } n = 7 \\
  S_7 &= \frac{2(4^7 - 1)}{4 - 1} \\
  S_7 &= \frac{2(16384 - 1)}{3} \\
  S_7 &= \frac{2 \times 16\,382}{3} \\
  S_7 &= 10922
\end{align*}
\booknote{$4^7 - 1 = 16383$; the book's ``$16\,382$'' is a misprint, though the
final answer $10922$ is correct.}
\end{example}

\begin{example}{Example 2.5}
Find the sum of the first 10 terms of $\frac{1}{2}, \frac{1}{4}, \frac{1}{8}, \ldots$

\solution
\begin{align*}
  S_n &= \frac{a(1 - r^n)}{(1 - r)},\ r \ne 1 \\
  a &= \frac{1}{2},\ r = \frac{1}{2} \text{ and } n = 10 \\
  S_{10} &= \frac{\frac{1}{2}\left(1 - \left(\frac{1}{2}\right)^{10}\right)}{\left(1 - \frac{1}{2}\right)} \\
  S_{10} &= \frac{1}{2}\left(1 - \frac{1}{1024}\right) \\
  S_{10} &= \frac{1024 - 1}{1024} \\
  S_{10} &= \frac{1023}{1024}
\end{align*}
\booknote{the line $S_{10} = \frac{1}{2}\left(1 - \frac{1}{1024}\right)$ is printed as in the
book; dividing by $\left(1 - \frac{1}{2}\right)$ cancels the $\frac{1}{2}$, which is why the
next line is $\frac{1024-1}{1024}$.}
\end{example}
